%% ====================================================================
\section{Objetivos}
\label{sec:objetivos}
%% ====================================================================

A partir del problema planteado en la introducción, esta sección formula el objetivo
general del TFM y lo desglosa en objetivos específicos, medibles y alcanzables. Los
objetivos específicos guían el desarrollo metodológico y la campaña experimental, y su
grado de cumplimiento se contrastará en la entrega final.

\subsection{Objetivo general}

Diseñar y evaluar, mediante simulación cuantitativa reproducible, un esquema de
coordinación distribuida local basado en consenso dinámico de media y dinámicas
poblacionales para la formación adaptativa de coaliciones en flotas multi-AGV que
transportan cooperativamente cargas heterogéneas con requisitos mínimos de
cardinalidad, integrando estimación de ocupación, asignación a cargas, formación de
coaliciones y navegación reactiva en un único bucle de control, caracterizado mediante
una escalera comparativa de cinco tratamientos (T1--T5) y contrastado con una
referencia centralizada replanificada bajo seis escenarios experimentales.

\subsection{Objetivos específicos}

\begin{enumerate}
  \item Formalizar el problema de formación de coaliciones multi-AGV como un
    sistema de asignación distribuida con requisitos mínimos de cardinalidad,
    estrategia de inactividad, comunicación local y navegación cinemática en un
    entorno 2D.

  \item Diseñar una función de payoff con componente sigmoide decreciente de
    cardinalidad y descuento espacial, distinguiendo entre el juego base con
    $\Psi\equiv1$ (donde se analiza la estructura de juego potencial y se caracteriza
    el intervalo $(\beta_{\min}, \beta_{\max})$ de estabilidad local del equilibrio
    mediante un análisis de Lyapunov) y el sistema implementado con $\Psi(d)$,
    tratado como una perturbación espacial validada por simulación.

  \item Proponer la tasa de revisión espacialmente modulada $\mu_i(t)$ y
    cuantificar su efecto sobre la frecuencia de reasignaciones y el tiempo de
    formación de coaliciones en los experimentos.

  \item Implementar las tres dinámicas acopladas (consenso distribuido,
    Smith en tiempo discreto y gradiente reactivo con evitación de obstáculos
    estáticos) en un paquete Python
    reproducible, con soporte para cinco tratamientos y seis escenarios con semillas
    controladas y registro completo de trayectorias.

  \item Validar cualitativamente el comportamiento de formación de coaliciones y el
    transporte cooperativo de la coalición con evitación reactiva de obstáculos en
    CoppeliaSim con robots Pioneer~3-DX.
\end{enumerate}
